(1) 阶跃驱动后常系数方程 \(\tau\dot P+P=\varepsilon_0\Delta\varepsilon E_0\)，用 \(P(0)=0\) 得 \(\boxed{P(t)=\varepsilon_0\Delta\varepsilon E_0(1-e^{-t/\tau})}\)（\(t\ge0\)）。对任意历史电场的稳态因果卷积为 \(P(t)=\varepsilon_0\Delta\varepsilon\int_{-\infty}^t e^{-(t-t')/\tau}E(t')dt'/\tau\)，核仅在过去 \(t'\le t\) 非零，体现物理因果性。瞬时位移部分为 \(\varepsilon_0\varepsilon_\infty E\)，不属于慢极化变量 \(P\)。

(2) 代入谐波形式 \(P_\omega=\varepsilon_0\Delta\varepsilon E_0/(1-i\omega\tau)\)，所以 \(\boxed{\varepsilon^*(\omega)=\varepsilon_\infty+\Delta\varepsilon/(1-i\omega\tau)}\)。有 \(\varepsilon'=\varepsilon_\infty+\Delta\varepsilon/[1+(\omega\tau)^2]\)、\(\varepsilon''=\Delta\varepsilon\,\omega\tau/[1+(\omega\tau)^2]\)。\(\omega\to0\) 时 \(\varepsilon'\to\varepsilon_\infty+\Delta\varepsilon,\varepsilon''\to0\)；\(\omega\to\infty\) 时 \(\varepsilon'\to\varepsilon_\infty,\varepsilon''\to0\)。采用 \(e^{-i\omega t}\) 时被动损耗虚部为正。

(3) 周期平均电场对慢极化做功率密度为 \(\langle E\dot P\rangle=(\omega\varepsilon_0\varepsilon''E_0^2)/2\)，乘周期 \(2\pi/\omega\) 得 \(\boxed{W_{\rm cycle}=\pi\varepsilon_0\varepsilon''E_0^2}\ge0\)。令 \(x=\omega\tau\)，\(\varepsilon''=\Delta\varepsilon x/(1+x^2)\)，导数 \(\Delta\varepsilon(1-x^2)/(1+x^2)^2\)，故最大值在 \(x=1\)，为 \(\boxed{\varepsilon''_{\max}=\Delta\varepsilon/2}\)。介质损耗是弛豫相位滞后造成的不可逆能量传递，不能把瞬时实介电响应误认为耗散。

(4) 将上问 \(\varepsilon''\) 代入色散和则，得到 \((2/\pi)\int_0^\infty\varepsilon''d\omega/\omega=(2\Delta\varepsilon/\pi)\int_0^\infty\tau\,d\omega/[1+(\omega\tau)^2]=\Delta\varepsilon\)。该积分正好等于 \(\varepsilon'(0)-\varepsilon'(\infty)\)，是单极点因果响应的Kramers–Kronig检查。若 \(\Delta\varepsilon\to0\)，慢极化与损耗消失；若 \(\tau\to0\) 且观察固定有限 \(\omega\)，慢极化近乎瞬时跟随、损耗趋零，但无限频积分和则的极限不可与积分交换。
